\documentclass[12pt,a4paper]{article}

\usepackage{german}
\usepackage[utf8]{inputenc}
\usepackage{amsfonts}
\usepackage{amsthm}
\usepackage{amssymb}
\usepackage{enumerate}

\usepackage{graphicx}

\oddsidemargin=0cm \headheight=-2.1cm
\textwidth=16.4cm \textheight=26cm

\begin{document}
\thispagestyle{empty}
\begin{figure}[h]
\includegraphics{it-smo_logo}
\end{figure}

\vspace{1cm}

\begin{center}
\Huge{\textbf{Turno preliminare SMO 2017}}\\[1.5cm]
\large{Lugano, Losanna, Zurigo - 14 gennaio 2017}\\[3.5cm]
\end{center}


Siete pregati di rispettare i seguenti punti.

\begin{itemize}
\item Le vostre soluzioni saranno fotocopiate e suddivise per esercizio. Scrivete quindi solo su un lato dei fogli e cominciate un foglio nuovo per ogni esercizio. Qualsiasi dettaglio che possa facilitare la correzione è benvenuto, provate per esempio a evidenziare i passaggi importanti e a marcare ciò che è suprefluo.

%Vos solutions seront copi�es et tri�es par exercice. N'�crivez
%donc que sur un c�t� des feuilles et commencez une nouvelle
%feuille pour chaque exercice. Tous les d�tails facilitant les
%corrections sont les bienvenus, essayez par exemple de mettre en
%�vidence les passages importants et de marquer ce qui est superflu.

\item Scrivete TUTTO ciò che trovate, anche delle osservazioni banali possono valere dei punti. Consegnate TUTTI i fogli leggibili, anche per affermazioni completamente sba\-glia\-te non togliamo mai dei punti.
%�crivez TOUT ce que vous trouvez, m�me des observations
%triviales peuvent valoir des points. Et rendez TOUTES les feuilles
%lisibles, m�me pour des assertions compl�tement fausses on n'enl�ve
%jamais de points.

\item Gli esercizi sono ordinati per difficoltà ma capita spesso di non vedere la soluzione di un esercizio semplice ma di risolverne perfettamente uno difficile. Non lasciatevi dunque impressionare dal numero di un esercizio. Tutti gli esercizi valgono lo stesso numero di punti. Avete abbastanza tempo per occuparvi di ogni esercizio e vi consigliamo vivamente di farlo.

%Les exercices sont tri�s par difficult� mais il arrive souvent
%que quelqu'un ne voit pas la solution d'un exercice simple tout en
%r�solvant parfaitement un des probl�mes difficiles. Ne vous laissez
%donc pas impressionner par le num�ro d'un exercice. Tous les
%exercices valent le m�me nombre de points. Vous avez assez de temps
%pour vous occuper de chaque exercice et on vous conseille vivement
%de le faire.

\item È possibile consegnare l'esame prima della fine e lasciare l'aula senza fare rumore.

%Il est possible de rendre l'examen avant la fin et de quitter
%la salle sans faire de bruit.

\item Per evitare di perdere qualcosa durante le correzioni mettete su ogni foglio il vostro nome (o almeno un'abbreviazione), il numero della pagina e il numero totale di pagine per quell'esercizio. Per esempio ``L. Euler, G1), p. 2/4''.

%Pour qu'on ne perde rien en corrigeant mettez � la fin de
%l'examen sur chaque page votre nom (ou au moins une abr�viation), le
%num�ro de la page et le nombre total de pages pour cet exercice. Par
%exemple "`A.Cauchy No 3 p. 2/4"'.

%\item 
%On va t'envoyer ta copie corrig�e la semaine suivant l'examen
%dans l'enveloppe ci-joint. S'il te pla�t �cris maintenant dessus ton
%adresse correcte et bien lisible.

\end{itemize}
\end{document}
